% !TeX program = xelatex
% Self-contained English source. Run XeLaTeX twice.
\documentclass[11pt,a4paper]{article}
\usepackage[margin=23mm,top=22mm,bottom=20mm,headheight=15pt]{geometry}
\usepackage{fontspec,amsmath,amssymb,booktabs,tabularx,amsthm,fancyhdr,xcolor,hyperref}
\IfFontExistsTF{Noto Serif}{\setmainfont{Noto Serif}}{}
\IfFontExistsTF{Noto Sans}{\setsansfont{Noto Sans}}{}
\definecolor{teal}{HTML}{155955}\definecolor{ink}{HTML}{1F3840}
\hypersetup{colorlinks=true,urlcolor=teal,linkcolor=teal,pdftitle={Two proof attempts in spectral realization and tiling},pdfauthor={Shaitetsu}}
\pagestyle{fancy}\fancyhf{}\renewcommand{\headrulewidth}{.25pt}
\fancyhead[L]{\sffamily\footnotesize Spectral realization and tiling}\fancyhead[R]{\sffamily\footnotesize Research note}
\fancyfoot[L]{\sffamily\footnotesize 15 November 2026}\fancyfoot[R]{\sffamily\small\thepage}
\setlength{\parindent}{0pt}\setlength{\parskip}{5pt}\setlength{\emergencystretch}{2em}
\newcommand{\kicker}[1]{\par{\sffamily\bfseries\color{teal}\small #1}\par}
\newcommand{\sect}[1]{\section*{#1}}
\newcommand{\R}{\mathbb R}\newcommand{\Z}{\mathbb Z}\newcommand{\C}{\mathbb C}
\newcommand{\ee}{\mathrm e}
\begin{document}
\pdfbookmark[0]{Results and scope}{intro}
\kicker{RESEARCH NOTE\quad /\quad 15 NOVEMBER 2026}
\vspace{3mm}
{\sffamily\bfseries\color{ink}\fontsize{29}{35}\selectfont Two proof attempts in spectral realization and tiling\par}
\vspace{2mm}
{\large Exact realization conditions for repeated eigenvalues and a sharp test for translation tiling on the real line.\par}
\vspace{4mm}
\kicker{I. WHICH EIGENVALUE LISTS CAN A NONNEGATIVE MATRIX REALIZE?}
A symmetric matrix with smallest eigenvalue $-t$ repeated $n-2$ times has a planar Gram representation after adding $tI$. We classify all such spectra in every dimension $n\ge3$. With a prescribed diagonal, the test is a subset-sum condition. With a free diagonal, only two scalar inequalities remain. We also classify a separate constrained problem for five vertices with one missing edge and every row sum equal to one. At smallest eigenvalue $-9/20$, the sharp lower bound for the second eigenvalue is $25/32$ and the endpoint is excluded.
\vspace{3mm}
\kicker{II. HOW MUCH FOURIER ORTHOGONALITY FORCES A TILE?}
Let $\Omega\subset\R$ be measurable with Lebesgue measure one. Among integers $d\ge2$, orthogonality of all frequencies in $d\Z$ and just one frequency outside it forces a translational tiling for every such $\Omega$ exactly when $d=2$ or $3$. For every $d\ge4$ there is a counterexample to that implication. For every $d$, the sharp number of consecutive harmonics is $\lfloor d/2\rfloor$. At that threshold the proof gives both an explicit tiling complement and a full exponential basis. With one fewer harmonic, a nontiling set exists for each $d$.
\vspace{3mm}
\kicker{SCOPE AND METHOD}
The two unrestricted problems remain unresolved by these arguments. The matrix proof uses geometry of a positive semidefinite shift and saturates the bounds by a construction. The Fourier proof uses periodization, Newton identities and exact finite tiling obstructions. The methodological starting points in the earlier study were VMRC cases 015 (the matrix problem) and 056, 068, 185 (character sums and finite packing). Related theorems and source links appear on the last page. No claim of historical priority or formal proof-assistant verification is made.
\clearpage
\pdfbookmark[0]{Rank-two shifts in arbitrary dimension}{ranktwo}
\kicker{I\quad /\quad SYMMETRIC NONNEGATIVE MATRICES}
\sect{An exact test in arbitrary dimension}
Fix $n\ge3$, $t>0$ and $\rho\ge b>-t$. Write
\[
 \alpha=\rho+t,\quad \beta=b+t,\quad T=\alpha+\beta,\quad \Delta=\alpha-\beta .
\]
\textbf{Theorem 1.} An entrywise nonnegative symmetric matrix with spectrum $(\rho,b,-t,\ldots,-t)$ exists exactly when
\[
             \boxed{\rho+b-(n-2)t\ge0,\qquad \rho+t\ge\lceil n/2\rceil t.}
\]
With a specified diagonal $d_i\ge0$, let $L_i=d_i+t$. The exact criterion is $\sum_i L_i=T$ and the existence of a subset $S$ with
\[
                     \beta\le\sum_{i\in S}L_i\le\alpha.
\]
\paragraph{Planar reduction.} The shifted matrix $A+tI$ is positive semidefinite of rank two. Its Gram vectors have squared lengths $L_i>0$ and all pairwise inner products are nonnegative. Orient one vector at angle zero and place all angles in $[-\pi/2,\pi/2]$. The two extreme angles differ by at most $\pi/2$, because their cosine is nonnegative. Rotate them into $[0,\pi/2]$. If the angles are $\theta_i$, the two eigenvalues of the $2\times2$ frame operator have sum $\sum_iL_i$ and gap
\[
        \left|\sum_i L_i\exp(2\mathrm i\theta_i)\right|.
\]
This follows by subtracting the two diagonal entries and doubling the off-diagonal entry. Its two eigenvalues are also the nonzero Gram eigenvalues.
\paragraph{Necessity and construction.} Write $\phi_i=2\theta_i\in[0,\pi]$. Keeping the other vectors fixed with sum $z$, the squared resultant differs from a constant by $2L_i\Re(\bar z\exp(\mathrm i\phi_i))$. Since $\Im z\ge0$, its minimum on the upper semicircle is attained at $0$ or $\pi$. Rounding each angle to an endpoint in turn cannot increase the resultant. Hence a signed partition of weights has imbalance at most $\Delta$, equivalently one part has weight in $[\beta,\alpha]$. Conversely let the two weights be $U,V$. Put the two groups at doubled angles $0,\phi$. The resultant varies continuously from $T$ to $|U-V|\le\Delta$ as $\phi$ runs from $0$ to $\pi$. Use an angle giving $\Delta$. The Gram matrix minus $tI$ has the desired diagonal, signs and spectrum.
\paragraph{Eliminating the diagonal.} Each part of the partition weighs at most $\alpha$. One part has at least $k=\lceil n/2\rceil$ members, giving $\alpha\ge kt$, and $T\ge nt$ follows by summing. Conversely, prescribe $k$ members for the first part and $n-k$ for the second. Set $U=\max(T/2,kt)$ and $V=T-U$. Then $U\le\alpha$, $V\ge(n-k)t$ and $V\le U$. Each part can distribute its weight with every $L_i\ge t$. The construction above applies.\hfill$\square$
\paragraph{A list separated by the second inequality.} For $(1,9/10,-3/5,-3/5,-3/5)$, the trace is $1/10$, but $1+3/5<3(3/5)$. Thus no such symmetric nonnegative matrix exists. The second inequality supplies information even for a positive-trace list.

\clearpage
\pdfbookmark[0]{The five-vertex stochastic slice}{stoch}
\kicker{I\quad /\quad PRESCRIBED SUPPORT AND ROW SUMS}
\section*{The stochastic five-vertex slice}
Fix $1>b>a$ and prescribe $(1,b,a,a,a)$ as the spectrum of a symmetric
entrywise nonnegative matrix with row sums one. We prescribe exactly one
missing off-diagonal edge, so $A_{12}=A_{21}=0$ and all remaining
off-diagonal entries must be strictly positive. Write
\[
 u=\frac{1-a}{5},\qquad R=\frac{b-a}{u},\qquad
 q=\max\left(0,-\frac{a+u}{u}\right),\qquad d=\sqrt q.
\]
The prescribed matrix exists exactly in the following ranges:
\[
\begin{array}{c|l}
 0\le q<1/2 & 2+4q\le R<5\\
 1/2\le q<2/3 & 12q-6+2/q<R<5\\
 q\ge2/3 & \text{no realization}.
\end{array}
\]

\paragraph{Coordinates.}
Since the Perron eigenspace is generated by the all-ones vector,
the spectral theorem gives $A=aI+u(J+xx^{\mathsf T})$ with $J$ the
all-ones matrix, $\sum x_i=0$ and $\sum x_i^2=R$. Conversely these
conditions give the prescribed row sums and spectrum. Diagonal
nonnegativity means $|x_i|\ge d$. At the missing edge we can take
$x_1=X>0$ and $x_2=-Y<0$, where $XY=1$. All three remaining coordinates
must lie strictly between $-Y$ and $X$. This is also sufficient for
strict positivity on the remaining edges: every other opposite-sign
product has absolute value strictly less than $XY$.

\paragraph{Lower bounds.}
Write $h=X-Y$. Then $X^2+Y^2=2+h^2$. For $d=0$ this already gives
$R\ge2$. Suppose $d>0$. Change the sign of $x$ and relabel the two
endpoints if necessary so that its positive sign group is the smaller.
If it consists only of $X$, then $h\ge3d$ and $R\ge2+12d^2$.
Here $Y\ge d$, $X=1/Y$, and $X-Y\ge3d$ imply $d^2\le1/4$.
In the other case the remaining coordinates are $s,-w_1,-w_2$, where
$d\le s<X$, $d\le w_i<Y$ and $h=w_1+w_2-s$. In particular $X,Y>d$.
If $h\ge d$, then $R\ge2+4d^2$. If $h<d$, then $s\ge2d-h$ and
\[
 R\ge 2+h^2+2d^2+(2d-h)^2
   =2+4d^2+2(h-d)^2 .
\]
For $q\ge1/2$ the one-positive case is impossible, and $Y>d$ gives
$h<1/d-d\le d$. Substitution yields the strict bound
$R>2+4d^2+2(1/d-2d)^2=12q-6+2/q$.
The positive total $X+s<2X<2/d$ must equal the negative total,
which is at least $3d$. Thus $q\ge2/3$ is impossible. Finally
$R<5$ follows from $b<1$.

\newpage
\section*{Construction throughout the admissible region}
It suffices to set $w_1=w_2=w$. Given $X>0$, put $Y=1/X$, $h=X-Y$
and $s=2w-h$. Then
\[
       x=(X,-Y,s,-w,-w),\qquad
       R(X,w)=2+h^2+(2w-h)^2+2w^2 .
\]
We impose
\[
 X>d,\quad Y>d,\qquad
 L(X):=\max\{d,(d+h)/2\}\le w
 < U(X):=\min\{Y,(2X-Y)/2\}.
\]
These conditions ensure $d\le s<X$ and $d\le w<Y$; they apply also
when $d=0$, in which case zero interior coordinates are permitted.

\paragraph{Why no values of $R$ are missed.}
The $X$-domain is defined by $X>d$, $Y>d$, $2X-Y>2d$ and
$3Y-X>d$: indeed these four inequalities are equivalent to the two
endpoint requirements and $L(X)<U(X)$. Both $X$ and $2X-1/X$ increase
and both $1/X$ and $3/X-X$ decrease. The domain is therefore an
interval. Over it the vertical sections $[L(X),U(X))$ are nonempty
and the lower endpoint is continuous. Any two feasible points can be
joined by moving down to that continuous graph and following it.
Consequently the image of $R(X,w)$ is an interval.

If $0\le q<1/2$, choose $X,Y$ satisfying $XY=1$, $X-Y=d$. Then
$Y>d$ exactly when $q<1/2$, and $w=s=d$ is admissible and attains
$R=2+4q$. This also covers $d=0$, with $X=Y=1$ and $s=w=0$.
If $1/2\le q<2/3$, let $Y$ tend to $d$ from above, with $X=1/Y$
and $w=d$. Here $h<1/d-d\le d$, so $s>d$. At the limiting point
$s<X$ follows from $3d<2/d$ and $X>d$ from $q<1$. Both inequalities
persist for $Y$ sufficiently close to $d$. The limiting value is
$2+4d^2+2(1/d-2d)^2$, never attained by the strict lower bound.

For either range take $X=\sqrt{3/2}$, $Y=\sqrt{2/3}$ and let $w$
increase to $Y$ from below. Because $2X=3Y$ and $d<Y$, these parameters
are admissible sufficiently close to the limit. Then $s$ tends to $X$
and $R$ tends to $2X^2+3Y^2=5$. The interval property now proves
sufficiency for every stated value. For example, with $a=-9/20$ the
criterion is exactly $25/32<b<1$.
\clearpage
\pdfbookmark[0]{A certified box of simple spectra}{simple}
\kicker{I\quad /\quad SPLITTING REPEATED EIGENVALUES}
\sect{An open box with row sums and support fixed}
The three equal eigenvalues in the five-vertex problem can be varied independently near the following point. Write
\[
 \lambda^0=(8347/10000,-9/20,-9/20,-9/20).
\]
For every real four-vector $\lambda$ with $\|\lambda-\lambda^0\|_\infty<10^{-4}$, the spectrum $(1,\lambda_1,\ldots,\lambda_4)$ is realized by a symmetric nonnegative matrix with row sums one and exactly one missing edge. No lower bound on the separation between the four eigenvalues is needed.
\paragraph{A base point with positive margins.}
Use $a=-9/20$, $u=29/100$, and
\[
 x=(5/4,-4/5,21/20,-3/4,-3/4)^\mathsf T,\qquad
 A_0=aI+u(J+xx^\mathsf T).
\]
The coordinates of $x$ sum to zero, $\|x\|^2=443/100$, and $a+u\|x\|^2=8347/10000$. The missing entry is $(A_0)_{12}=0$, the smallest diagonal is $1/320$, and the smallest positive off-diagonal is $29/1600$. Denote the normalized all-ones vector by $e_0$ and $x/\|x\|$ by $e_1$. Complete these to an orthonormal basis $e_0,\ldots,e_4$. Set
\[
 v=(0,9,1,-10,0)^\mathsf T,\quad
 K=(vx^\mathsf T-xv^\mathsf T)/\|x\|^2,\quad Q(s)=\exp(sK).
\]
Here $v\perp x,e_0$, $Kx=v$, $Ke_0=0$, and the operator norm $\|K\|=\|v\|/\|x\|<7$. The matrices
\[
 A(s,\lambda)=Q(s)\left(e_0e_0^\mathsf T+\sum_{j=1}^4\lambda_j e_je_j^\mathsf T\right)Q(s)^\mathsf T
\]
have the desired row sums and eigenvalues. Only the entries need control.
\paragraph{Solving the missing-edge equation.}
Write $F(s)=A(s,\lambda^0)_{12}$. We have $F(0)=0$ and
$F'(0)=u(v_1x_2+x_1v_2)=261/80$. Differentiating the orthogonal conjugation twice gives $|F''(s)|\le4\|K\|^2\|A_0\|<196$, where $\|A_0\|=1$. For $\delta=10^{-4}$, Taylor's theorem gives opposite signs at $s=\pm\delta$, each of magnitude at least $(261/80)\delta-98\delta^2>\delta$. Perturbing $\lambda$ within the box changes each entry at a fixed $s$ by less than $\delta$. Thus the intermediate value theorem supplies $|s|<\delta$ with $A(s,\lambda)_{12}=0$. Meanwhile $\|A(s,\lambda)-A_0\|<\delta+14\delta<1/320$. Every required positive entry stays positive.\hfill$\square$
\paragraph{A specific simple spectrum.}
An example in the box is
\[ (1,8347/10000,-9/20+1/20000,-9/20,-9/20-1/20000). \]
 This establishes a four-dimensional neighborhood. It leaves the boundary of the full four-dimensional realizability region undetermined.
\clearpage
\pdfbookmark[0]{Orthogonality forces a tiling and a basis}{harmonics}
\kicker{II\quad /\quad FOURIER ANALYSIS ON THE REAL LINE}
\sect{Completing an orthogonal family}
Write $e_\xi(y)=\exp(2\pi\mathrm i\xi y)$. A spectrum is a set of frequencies whose exponentials form an orthonormal basis for $L^2(\Omega)$, with $|\Omega|=1$. A tiling by $\mathcal T$ means $\sum_{s\in\mathcal T}1_\Omega(y-s)=1$ almost everywhere.
\paragraph{Theorem 2.} Fix $d\ge2$, $h=\lfloor d/2\rfloor$. Let $\Omega\subset\R$ be Lebesgue measurable and of measure one. Suppose that $d\Z$ is orthogonal and that $e_{j\lambda}$ is orthogonal to every $e_{kd}$ for $1\le j\le h$, $k\in\Z$. In particular $j\lambda\notin d\Z$. Then $\zeta=\exp(2\pi\mathrm i\lambda/d)$ has finite order $dg$ for a positive integer $g$. A spectrum and a tiling set for $\Omega$ are
\[
 \Lambda'=\bigcup_{j=0}^{d-1}(j\lambda+d\Z),
 \qquad \mathcal T=\tfrac1d\{0,1,\ldots,g-1\}+g\Z .
\]
\paragraph{Count and phases on each fiber.} For $0\le x<1/d$ set $B_x=\{n\in\Z:x+n/d\in\Omega\}$. Tonelli gives $\int_0^{1/d}|B_x|\,dx=1$, so all periodizations below are integrable. The nonconstant Fourier coefficients of $|B_x|$ vanish by orthogonality on $d\Z$. Uniqueness of Fourier series gives $|B_x|=d$ almost everywhere. For each $j=1,\ldots,h$, all Fourier coefficients of
\[
 x\longmapsto e_{j\lambda}(x)\!\sum_{n\in B_x}\zeta^{jn}
\]
vanish by the cross-orthogonality assumption, so each sum vanishes almost everywhere. Take the common full-measure set. The $d$ unit numbers $z_n=\zeta^n$ on any such fiber have power sums zero in degrees $1,\ldots,h$. Newton's identities successively give vanishing elementary symmetric functions $E_1,\ldots,E_h$. The identity $E_{d-j}=E_d\overline{E_j}$ covers all intermediate degrees. Consequently the monic polynomial with the $z_n$ as roots is $z^d-c_x$ with $|c_x|=1$. Its roots are distinct and form a rotated regular $d$-gon.
\paragraph{One complement for all fibers.} Two phases on a good fiber have a primitive $d$th root as their ratio. Thus some nonzero power of $\zeta$ has order $d$, so $\zeta$ has a finite order $M$ divisible by $d$. Write $M=dg$. Under $n\mapsto\zeta^n$, the preimage in $\Z_M$ of the $d$th roots of unity is $g\Z_M$. Every good $B_x$ therefore represents exactly the elements of one coset of $g\Z_M$ modulo $M$. In $\Z_M$ its common tiling complement is $\{0,\ldots,g-1\}$. Lifting the complement periodically to $\Z$ and dividing translations by $d$ gives the asserted tiling of $\R$.
\paragraph{Completeness.} On each good fiber the $d$ columns $(z_n^j)_{n\in B_x}$, $0\le j<d$, are orthogonal of squared norm $d$ by the regular polygon calculation. Integrating first on fibers gives orthogonality within $\Lambda'$. If $f$ is orthogonal to $\Lambda'$, for each $j$ ordinary Fourier uniqueness on $[0,1/d)$ gives $\sum_{n\in B_x}f(x+n/d)\overline{z_n^j}=0$ almost everywhere, after removing the nonzero base phase. Invertibility of the fiber matrix gives $f=0$.\hfill$\square$
\clearpage
\pdfbookmark[0]{Sharpness of one extra frequency}{sharp}
\kicker{II\quad /\quad WHEN ONE ADDITIONAL FREQUENCY SUFFICES}
\section*{One additional frequency, including the first failure}
For each $d\ge2$, let $h=\lfloor d/2\rfloor$. Theorem 2 supplies an upper bound of $h$ extra harmonics. On the next page we prove it is optimal for each $d$ even when the first omitted harmonic is a new frequency. For $d=2,3$, the single extra frequency required by the theorem is automatic if $\Omega$ has a full spectrum containing a translate of $d\Z$. Indeed after shifting the frequencies back, $d\Z$ alone spans only functions constant on each $d$-point fiber. For example, assigning $+1$ and $-1$ at the smallest two indices of each good fiber and zero elsewhere gives a nonzero measurable function orthogonal to the lattice.
\paragraph{A two-interval counterexample at $d=4$.}
Set
\[
 B=\{0,1,3,4\},\qquad \Omega_*=(B+[0,1))/4=[0,1/2)\cup[3/4,5/4).
\]
Then $4\Z\cup\{2\}$ is orthogonal. For differences in $4\Z\setminus\{0\}$ integrate over each unit digit cell before scaling. For differences in $2+4\Z$ the digit sum is $1-1-1+1=0$. In fact $2\Z$ is orthogonal.
\paragraph{Why no real translates can tile.}
For any finite digit set $F$ and any translational tiling of $F+[0,1)$ by $S\subset\R$, choose $x\in(0,1)$ such that every $x+n$, $n\in\Z$, is covered exactly once. Such $x$ exist since there are countably many tests. The translate $s$ covers $x+n$ precisely when $n\in F-\lfloor x-s\rfloor$. Hence $F$ tiles $\Z$, with unique representations. Repeated translation indices would cover an integer twice. For $F=B$, the positive differences of $B$ contain $1,2,3,4$, so any two translation indices would have to be at least $5$ apart. A length-$N$ interval could then contain at most $N/5+O(1)$ indices, covering at most $4N/5+O(1)$ of its integers, after accounting for translates meeting the endpoints. This contradiction proves the claimed failure.
\paragraph{A counting obstruction for the general construction.}
Write $P_F(X)=\sum_{a\in F}X^a$ for nonnegative integer digits with minimum zero, translating the digits first if necessary. If $F$ tiles $\Z$, its complement indicator satisfies $y_n=1-\sum_{a\in F,\,a>0}y_{n-a}$. This finite-state recurrence is reversible using $\max F$, hence has a period $M$. Let $C$ be the complement digits in one period. Then $|F||C|=M$. For a prime $p$ and $1\le j\le v_p(M)$, the mask identity modulo $X^M-1$ forces $\Phi_{p^j}$ to divide $P_F$ or $P_C$. At most $v_p(|C|)$ of these distinct irreducibles divide $P_C$, since each takes the value $p$ at one. Thus at least $v_p(|F|)$ must divide $P_F$. This is the necessary prime-power count in the Coven--Meyerowitz tiling condition.
\clearpage
\pdfbookmark[0]{Optimality for every integer spacing}{optimal}
\kicker{II\quad /\quad MATCHING LOWER BOUNDS}
\section*{The exact number of consecutive harmonic tests}
For every $d\ge2$ there exists a nontile of measure one and a $\lambda$ for which the first $h-1$ harmonic tests hold, $h=\lfloor d/2\rfloor$. We can keep $h\lambda$ outside $d\Z$ as well. Thus the optimality refers to this specific sequence of tests, with no restriction on other possible ways of recognizing a tile.
\paragraph{Digit roots.}
For $d$ distinct nonnegative digits $F$, the set $(F+[0,1))/d$ has measure one and $d\Z$ is orthogonal there. The cross-test for $j\lambda$ holds whenever $P_F(\exp(2\pi\mathrm i j\lambda/d))=0$. To exclude a factor $\Phi_{p^s}$, use residue counts. Its divisibility requires the $p$ counts in the lifts of every class modulo $p^{s-1}$ to be equal. Indeed the polynomial of residue counts modulo $p^s$ has degree less than $p^s$. If divisible by $\Phi_{p^s}=\sum_{i=0}^{p-1}X^{ip^{s-1}}$, its quotient has degree less than $p^{s-1}$, giving the asserted equality. Monomial translations preserve divisibility by these polynomials.
\paragraph{Odd $d=2h+1$.}
Let $L=h(h+1)$, choose $p\mid d$ prime and note $\gcd(p,L)=1$. In $\Z/L\Z$ take the following multiset of $d$ residues:
\[
   \{k(h+1):0\le k<h\}\ \sqcup\ \{kh:0\le k<h+1\}.
\]
Use the Chinese remainder theorem to lift the entries to distinct nonnegative digits, all $0\bmod p$ except one $1\bmod p$. Adding multiples of $pL$ separates repeated digits. Take $\lambda=d/L$. The two regular polygons of sizes $h$ and $h+1$ each have vanishing power sums for $1\le j<h$. The $h$th sum equals $h$, and $h<L$. The modulo-$p$ counts exclude $\Phi_p$. For $s\ge2$ the necessary lift-count equality for $\Phi_{p^s}$ would make every residue class modulo $p$ have size divisible by $p$. The singleton in class one excludes every such factor. Since $p\mid d$, the counting obstruction proves nontiling.
\paragraph{Even $d=2h\ge4$.}
Write $h=2^v u$ with $u$ odd. Put $L=3h$ and take two $h$-gons with residues $3k$ and $1+3k$ modulo $L$ for $0\le k<h$. Each class modulo $2^v$ has $2u$ entries. Split those entries between its two lifts modulo $2^{v+1}$. If $u>1$, use counts $1,2u-1$ in every class. If $u=1$, use $0,2$ in one class and $1,1$ in each other class. Here $v\ge1$ in the second case, so another class exists. These lifts are compatible since $v_2(L)=v$. Adding multiples of $2L$ produces distinct digits. Unequal counts exclude $\Phi_{2^{v+1}}$, and an odd count modulo $2^{v+1}$ excludes every higher two-power factor. At most $v$ factors remain although $v_2(d)=v+1$. Taking $\lambda=d/L$, the first $h-1$ sums vanish and the $h$th is $h(1+\exp(2\pi\mathrm i/3))\ne0$. The tested harmonics are mutually orthogonal as their differences are among the same orders. This also applies in the odd case. For $d=2$ the example $\Omega_*$ on the preceding page already admits $2\Z$. The weaker list of tests is empty and $\lambda=1$ is outside $2\Z$.\hfill$\square$
\clearpage
\pdfbookmark[0]{What remains and references}{refs}
\kicker{RESEARCH BOUNDARIES AND REFERENCES}
\sect{What the two arguments leave open}
\paragraph{Symmetric matrices.} The rank-two shift allowed all Gram vectors to lie in a quadrant and reduced the spectrum to an angular resultant. This reduction can fail immediately in rank three. Let $C$ be the adjacency matrix of the five-cycle and $\varphi=(1+\sqrt5)/2$. The nonnegative matrix $\varphi I+C$ is positive semidefinite with spectrum $(\varphi+2,\sqrt5,\sqrt5,0,0)$ by diagonalizing the cyclic shift. It has no entrywise nonnegative Gram factor, even with extra columns. Indeed each column of such a factor would have support in a clique of the cycle, of size at most two. The squared norm of such a column is at least twice its contribution to the five undirected edge entries. Summing would force the trace to be at least $10$, whereas it is $5\varphi<10$. Thus general configurations of pairwise nonobtuse vectors in rank three cannot all be given nonnegative coordinates. The small open box above provides examples in higher shifted rank but does not classify that rank.
\paragraph{Real-line tiling.} An arbitrary spectral set is supplied with a complete basis. The single-frequency counterexamples need only admit the specified orthogonal family. They say where incomplete Fourier data cease to force tiling. To prove the unrestricted forward implication one still needs a common tiling complement for the fibers of an arbitrary spectrum. The reverse implication, from a tiling to a spectrum, also remains outside these arguments.
\paragraph{Related work.} Dutkay and Haussermann, Theorem 1.26, already derived the forward implication for a complete spectrum of period $2,3,4$ or $5$. Their period-three proof uses the equilateral-triangle argument. The fiber methods here are adaptations of established Fourier methods. The harmonic test also proves completion to a full periodic basis directly from its incomplete input. The mask obstruction used in the sharpness proof is the necessary part of Coven and Meyerowitz's condition (T1).
\begin{itemize}\setlength\itemsep{2pt}
\item T. J. Laffey, \href{https://www.math.ualberta.ca/~vtroitsky/positivity2017/talks/Laffey.pdf}{On characterizing the spectra of nonnegative matrices} (2017).
\item D. E. Dutkay and P. E. T. Jorgensen, \href{https://arxiv.org/abs/1209.0666}{On the universal tiling conjecture in dimension one}.
\item D. E. Dutkay and J. Haussermann, \href{https://arxiv.org/abs/1210.8340}{Tiling properties of spectra of measures}, Math. Z. 280 (2015).
\item E. M. Coven and A. Meyerowitz, \href{https://arxiv.org/abs/math/9802122}{Tiling the integers with translates of one finite set}, J. Algebra 212 (1999).
\item \href{https://github.com/QEEDDD/Verifiable-Mathematical-Reasoning-Corpus/tree/5887cdea4898f0a3aa5bf7a8e0829697bd2b0f4b}{Verifiable Mathematical Reasoning Corpus}, snapshot 5887cde.
\end{itemize}
\end{document}
